\medskip\noindent\textbf{Supplementary Note 6. Ranking and uncertainty at matched batch removal.}\par\smallskip
\begingroup\small\setlength{\parskip}{3pt}
The six penalties are ranked by biology at the first crossing of $b^*$ (Supplementary Note 5). Of 16 task--decoder pairs, 13 are ranked. Supplementary Table S6 gives readings, crossings and uncertainty, with low-signal pairs separately. Of 78 ranked curves, 46 cross more than once.

Each reading is repeated with one seed omitted from both the penalty and plain scVI, holding $b^*$ fixed. For the $m$ available readings $x_{(i)}$, with mean $\bar x$, the jackknife standard error is
\begin{equation}
\mathrm{SE}=\Big(\frac{m-1}{m}\sum_{i=1}^{m}\big(x_{(i)}-\bar x\big)^2\Big)^{1/2},
\label{eq:s6-loso}
\end{equation}
It is reported when at least two repeats reach $b^*$. All 78 ranked readings had $m=3$. Each repeat also ranks all penalties. The rank range spans these repeat ranks and excludes the full-seed rank. It has the same availability rule as the standard error.

Rank 1 denotes highest biology. Readings within $10^{-9}$ tie and share their average rank. If $f$ penalties do not reach $b^*$, each receives the mean of places $7-f$ to 6. The six ranks therefore sum to 21. No full-seed sweep readings tied or failed to reach the target.

Mean ranks average ranked real tasks separately by decoder: five batch-conditioned and six unconditioned tasks. Simulations are ranked individually and never pooled with real tasks. Related tasks remain separate in these means. Human--mouse immune contains all 33,506 immune cells, and ATAC large contains all 11,270 ATAC small cells. Thus the six real tasks represent four families, two weighted twice (Supplementary Table S1).

Sensitivity analyses use per-seed curves and average per-seed ranks, omit trajectory conservation from immune biology scores, or remove strengths containing failed fits. The last reproduces the primary ranking because no fit failed (Supplementary Figure S2). Scaling analyses recompute scores, $b^*$ and ranks using min--max scaling, z-scores or a composite without silhouette metrics (Supplementary Figure S5). Scaling uses every scored fit of a task across decoders, penalties and baselines. A further variant omits each method family when setting minima and maxima. Min--max scaling and z-scores change the leading mean rank.

These rankings are descriptive. Interpolation depends on strength spacing, and nearby readings can exchange ranks on new seeds. Supplementary Note 10 discusses uncertainty limitations.
\par\endgroup
